\begin{abstract}
Η παρούσα διπλωματική εργασία παρουσιάζει τον \textbf{Πάγκο Εργασίας Εκπομπών} (\en{Emissions Workbench}), μια πλατφόρμα λογισμικού παραγωγικής κλίμακας που σχεδίασε και υλοποίησε η ομάδα \en{Energy Transition (ET) Platform} της A.P. Møller-Mærsk A/S. Στόχος της πλατφόρμας είναι η μέτρηση, πιστοποίηση και βελτιστοποίηση των εκπομπών αερίων θερμοκηπίου (\en{GHG}) από θαλάσσιες μεταφορές, στο πλαίσιο του στόχου \en{Net Zero 2040} της Maersk και των ρυθμιστικών υποχρεώσεων του κλάδου (CSRD, \en{EU ETS Maritime}, \en{FuelEU Maritime}).

Η εργασία αναλύει τέσσερις αλληλοσυμπληρούμενες διαστάσεις: το επιχειρησιακό και ρυθμιστικό πλαίσιο, το θεωρητικό υπόβαθρο λογιστικής εκπομπών (\en{GHG Protocol}, \en{Scope 1/2/3}, \en{ISCC}, \en{GLEC}), την τεχνική αρχιτεκτονική της πλατφόρμας και τη μεθοδολογία ανάπτυξης της ομάδας.

Τεχνικά, το σύστημα βασίζεται στην πλατφόρμα \en{BEAM/Elixir} με \en{Phoenix LiveView} για πραγματικού χρόνου διεπαφές, \en{Apache Kafka} για \en{event-driven} επεξεργασία δεδομένων, \en{event sourcing} για αδιάσπαστη ιχνηλασιμότητα της λογιστικής πράσινου καυσίμου (\en{Energy Bank}), \en{PostgreSQL/Ecto}, \en{Dremio} για αναλυτικά δεδομένα και \en{Microsoft Azure/Kubernetes} για αναπτύξεις νέφους. Η ομάδα ανάπτυξης εφαρμόζει πρακτικές \en{Extreme Programming}: \en{pair programming}, \en{test-driven development}, \en{continuous integration/delivery} και \en{vertical ownership}.

Επιχειρησιακά αποτελέσματα: πάνω από 1.300 ενεργοί χρήστες, παράδοση \en{ECO} προϊόντων ύψους 98.000 \en{FFE} αξίας \$31 εκατομμυρίων \en{USD} κατά το 2024, περίπου 32 αναπτύξεις σε παραγωγή ημερησίως και άνω των 720 \en{database migrations} χωρίς διακοπή υπηρεσίας.

\begin{keywords}
Εκπομπές αερίων θερμοκηπίου, CO\textsubscript{2}, Ναυτιλία, Maersk, Elixir, BEAM, Phoenix LiveView, Event Sourcing, Extreme Programming, Net Zero
\end{keywords}
\end{abstract}


\begin{abstracteng}
\begin{english}
This thesis presents the \textbf{Emissions Workbench}, a production-scale software platform designed and implemented by the Energy Transition (ET) Platform team at A.P. Møller-Mærsk A/S. The platform enables accurate measurement, certification, and optimization of greenhouse gas (GHG) emissions from maritime shipping operations, in support of Maersk's Net Zero 2040 commitment and the evolving regulatory landscape (CSRD, EU ETS Maritime, FuelEU Maritime).

The thesis examines four complementary dimensions: the business and regulatory context; the theoretical framework for GHG accounting (GHG Protocol, Scope 1/2/3, ISCC, GLEC); the technical architecture of the platform; and the software development methodology of the team.

The system is built on the BEAM/Elixir runtime with Phoenix LiveView for real-time web interfaces, Apache Kafka for event-driven data ingestion, event sourcing for the auditable green fuel accounting ledger (Energy Bank), PostgreSQL/Ecto for persistence, Dremio for analytics, and Microsoft Azure/Kubernetes for cloud deployment. The development team applies Extreme Programming practices: pair programming, test-driven development, continuous integration/delivery, and vertical ownership.

Business outcomes include over 1,300 active users, ECO Delivery volumes of 98,000 FFE worth \$31 million USD in 2024, approximately 32 production deployments per day, and over 720 zero-downtime database migrations.
\end{english}

\begin{keywordseng}
\begin{english}
Greenhouse gas emissions, CO\textsubscript{2}, Maritime shipping, Maersk, Elixir, BEAM, Phoenix LiveView, Event Sourcing, Extreme Programming, Net Zero
\end{english}
\end{keywordseng}
\end{abstracteng}
